\par\medskip
\begin{center}\textls[100]{e) Die letzten Tagebuchnotizen von 1800. Asymptotisches.}\\
\textls[100]{Mittelpunktsgleichung. Klassenanzahl.}\end{center}
\noindent\hspace*{\parindent}In das Jahr 1800 fallen noch einige auf analytische Gegenstände bezügliche
Tagebuchnotizen, mit deren Besprechung wir diesen Abschnitt beschliessen.
Die Aufzeichnung 112. vom 12. Juni berichtet von der »Erfindung eines ganz
neuen numerisch-exponentiellen Kalkuls«, womit die Methoden gemeint sind,
die in der Werke III, S. 426—431\footnote{Das Werke III, S. 431 mit den Worten »ecce iam computum pro \(e^{\frac12\pi}\)« beginnende, bis S. 432, stammt aus der Scheda Aa, ist also wesentlich älter und befolgt viel primitivere Methoden; vergl. oben S. 38 Fussnote \textsuperscript{54)}.} abgedruckten \textit{Sammlung von Rechnungen,
vornehmlich solchen, bei denen von meinen Methoden, die Faktoren grosser Zahlen
zu finden, und von den Wolframschen Logarithmentafeln Gebrauch gemacht ist}
angewandt werden\footnote{Diese Sammlung ist auf losen Blättern aufgezeichnet (Fh, 2).}. Die \textsc{Wolframschen} Logarithmentafeln geben die natürlichen
Logarithmen ganzer Zahlen auf 48 Dezimalstellen und sind in der
\textsc{Schulzeschen} Tafelsammlung (vergl. oben S. 11, \textsuperscript{14)}) enthalten. Der neue
Kalkul besteht nun darin, dass ein Vielfaches der zu berechnenden Exponentialgrösse,
wie z. B. der bei den lemniskatischen Funktionen auftretenden \(e^{-\pi}\),
\(e^{-\frac14\pi}\), durch eine ganze Zahl approximiert, und diese dann nach den Methoden
der \textit{Disqu. arithm.}\footnote{Diese werden in der Handschrift noch ohne Artikelnummer zitiert; es handelt sich um den Art. 329, Werke I, S. 401; vergl. dazu \textsc{Maennchen}, Abh. 6 dieses Bandes, Abschnitt VI, S. 27.} in Primfaktoren zerlegt wird; da nun die \textsc{Wolframschen}
Tafeln die Logarithmen sehr genau geben, so handelt es sich dann nur noch
um die Bestimmung von \(e^\delta\), wo \(\delta\) eine sehr kleine, sehr genau bekannte
Zahl ist.

Die Tagebuchnotiz 113. vom 25. Oktober 1800 berichtet über die »früher
vergebens gesuchte Lösung eines Problems der Wahrscheinlichkeitsrechnung
über Kettenbrüche«. \textsc{Gauss} hat in einem Briefe vom 30. Januar 1812 (Werke
X, 1, S. 371, 372) \textsc{Laplace} von diesem Problem Mitteilung gemacht, das er
»vor zwölf Jahren behandelt, aber nicht in befriedigender Weise gelöst« habe.
Es handelt sich dabei um die Bestimmung einer Funktion \(P(n,x)\), die der\par
\clearpage\SourceLeaf{200}\ArithmeticHead{115}{R}{Ve. ASYMPTOTISCHES. MITTELPUNKTSGLEICHUNG.}
\noindent Funktionalgleichung\footnote{Vergl. die Bemerkungen zu der Tagebuchnotiz 113. Werke X, 1, S. 552—556.}
\[
P(n+1,x)=\sum_{\nu=1}^{\infty}\left\{P\left(n,\frac1\nu\right)-P\left(n,\frac1{\nu+x}\right)\right\}
\]
genügt, während \(P(0,x)=x\), und \(P(1,x)\) mit der inexplikablen Funktion
\(1+\frac12+\cdot\cdot\cdot+\frac1x\) zusammenhängt; in \textsc{Gauss} späterer Bezeichnung (1812, \textit{Disqu.
circa seriem}, Werke III, S. 154) ist
\[
P(1,x)=\Psi(x)-\Psi(0).
\]
\textsc{Gauss} sagt, »für grössere Werte von \(n\) scheine die Funktion sich nicht behandeln
zu lassen, dagegen habe er durch sehr einfache Überlegungen gefunden,
dass für \(n\) gleich Unendlich
\[
P(n,x)=\frac{\log(1+x)}{\log2}
\]
sei«. Leider hat sich kein Fingerzeig dafür ergeben, wie \textsc{Gauss} diesen asymptotischen
Wert bestimmt haben mag. In dieser Hinsicht glücklicher sind wir
in bezug auf ein anderes Problem asymptotischen Charakters, dessen Behandlung
wohl auch in das Jahr 1800 fallen dürfte. \textsc{Gauss} schreibt nämlich am
5. Februar 1850 an \textsc{Schumacher} (Werke X, 1, S. 433): »Von meiner Methode,
den Grad der Konvergenz der nach Kosinus und Sinus der Vielfachen eines
Winkels fortschreitenden, eine beliebige periodische Funktion ausdrückenden
Reihe zu bestimmen, habe ich noch eine numerische Rechnung, welche sich
auf ein Beispiel der Mittelpunktsgleichung bezieht, aufgefunden, welches Blatt
wohl \(50\pm\) Jahre alt sein mag. Die Methode leistet aber viel mehr, als bloss
einen genäherten Ausdruck für ein sehr weit vom Anfang entferntes Glied
zu finden, sie ist auch geeignet, alle Glieder bis zum Anfang selbst hin
numerisch zu berechnen, und zwar mit aller zu wünschenden Schärfe.« Dieses
Blatt mit noch einigen andern, die sicher aus späterer Zeit (etwa 1805, vergl.
die historischen Feststellungen Werke X, 1, S. 443) stammen, hat sich im
Nachlass gefunden, ist Werke X, 1, S. 420—428 abgedruckt und ebenda
S. 438—442 erklärt. \textsc{Gauss} gibt daselbst für die Koeffizienten der Entwicklung
der Mittelpunktsgleichung \(v-M\) (wo \(v\) die wahre, \(M\) die mittlere Anomalie
bedeutet) nach den Sinus der Vielfachen von \(M\),\par
\HistoricalFoot{}{15\textsuperscript{*}}
\clearpage\SourceLeaf{201}\ArithmeticHead{116}{L}{L. SCHLESINGER, ÜBER GAUSS’ FUNKTIONENTHEORETISCHE ARBEITEN.}
\[
v-M=\sum_{n=1}^{\infty}\frac1n C_n\sin nM,
\]
\noindent den für grosse Werte von \(n\) geltenden asymptotischen Ausdruck
\[
C_n=\bigl(\operatorname{tg}\tfrac12\varphi\,e^{\cos\varphi}\bigr)^n
\left(1+\frac43\frac{1.}{\sqrt{2n\pi\cos^3\varphi}}\right),
\]
wo \(\sin\varphi\) die Exzentrizität ist. Dieser Ausdruck stimmt vollständig mit dem
überein, den \textsc{Jacobi} 1849\footnote{\textsc{Jacobis} Werke VII, S. 188, 237.}, einen früher von \textsc{Carlini} gegebenen verbessernd,
aufgestellt hat, aber die Herleitung von \textsc{Gauss} ist unvergleichlich kürzer als
die \textsc{Jacobis}, und bedient sich ganz elementarer Methoden. \textsc{Gauss} erreicht
sein Ziel wesentlich durch Anwendung der imaginären Transformation
\[
E=i\log\operatorname{cotg}\tfrac12\varphi+\varepsilon
\]
auf die exzentrische Anomalie \(E\); dabei ist es interessant festzustellen, dass
diese Transformation, wenn man sie auf den Integralausdruck
\[
C_n=\frac1\pi\int_{-\pi}^{+\pi}\frac{dv}{dM}\cdot e^{inM}dM
\]
jener Koeffizienten anwendet, sofort auf Integration im komplexen Gebiet
führt\footnote{Siehe Werke X, 1, S. 446.}, dass \textsc{Gauss} aber diesen Weg nicht eingeschlagen hat, obwohl ihm
sicher die Methode bekannt war, die \textsc{Laplace} (1781—1786) zur Bestimmung
solcher Integrale für grosse Werte der ganzen Zahl \(n\) angegeben hat. Er
hat vielmehr mit einer Reihendarstellung der \(C_n\) operiert.

In gewissem Sinne gehören in diesen Gedankenkreis auch die Tagebuchnotizen
114., 115. vom 30. November und 3. Dezember 1800, die die Entdeckung
von vier verschiedenen Methoden zur Bestimmung der Klassenanzahl
binärer quadratischer Formen anzeigen. — \textsc{Gauss} hat 1834 einen ersten Teil
seiner Untersuchungen \textit{Über den Zusammenhang der Klassenanzahl binärer quadratischer
Formen mit dem Wert ihres Determinanten} der Gesellschaft der Wissenschaften
vorlegen wollen; die Handschrift ist aus dem Nachlass Werke II,
S. 268 ff. abgedruckt und von \textsc{Dedekind} ausführlich erläutert worden\footnote{Vergl. auch \textsc{Bachmann}, Abhandlung 1 dieses Bandes, Artikel 27, S. 67.}. \textsc{Gauss}
beginnt den Art. 1 mit den Worten: »Schon 33 Jahre sind verflossen, seit\par
\clearpage\SourceLeaf{202}\ArithmeticHead{117}{R}{VIa. VON 1801 BIS ZUR BERUFUNG NACH GÖTTINGEN.}
\noindent wir die Prinzipien des wundersamen Zusammenhangs entdeckten, dem diese
Abhandlung gewidmet ist«, womit die Beziehung zu unseren Tagebuchnotizen
hergestellt erscheint\footnote{Vergl. auch die auf die Klassenanzahl bezügliche Stelle in dem Briefe an \textsc{Dirichlet} vom 2. November 1838, Werke XII, S. 310.}. Diese Prinzipien beruhen aber, wie \textsc{Gauss} im Art. 2
hervorhebt, wesentlich auf dem geometrischen Satze, dass das Verhältnis der
Anzahl der im Innern einer geschlossenen Kurve in der komplexen Zahlenebene
gelegenen Gitterpunkte mit ganzzahligen Koordinaten zu der eingeschlossenen
Fläche sich asymptotisch dem Wert 1 nähert, wenn die Kurve
gleichmässig ins Unendliche erweitert wird, also auf einem asymptotischen
Gesetz.
